
%! Author = lili
%! Date = 5/7/26

\newglossaryentry{frage}{
    name={test},
    description={todo},
    type=fragen
}

\newglossaryentry{dogma-arten}{
    name={Welche Arten des Informationstransfers gibt es laut dem zentralen Dogma?},
    description={DNA zu DNA (\gls{replikation}); DNA zu RNA (\gls{transkription}); RNA zu Protein (\gls{translation}); RNA zu DNA (\gls{rev-transkription})},
    type=fragen
}

\newglossaryentry{komplexitätmessen}{
    name={testWie kann man Genomkomplexität messen?},
    description={todo},
    type=fragen
}

\newglossaryentry{wieintrons}{
    name={Wie sind Gene mit Introns entstanden?},
    description={todo},
    type=fragen
}

\newglossaryentry{wiecmkarte}{
    name={Wie macht man eine genetische Karte?},
    description={Man kreuzt Organismen mit unterschiedlichen Merkmalen und analysiert die Nachkommen. Aus dem Anteil rekombinierter Nachkommen berechnet man die Rekombinationsfrequenz. Gene, die selten getrennt werden, liegen nahe beieinander; Gene mit hoher Rekombinationsrate liegen weiter auseinander. Durch den Vergleich vieler Genpaare bestimmt man schließlich die Reihenfolge der Gene und ihre relativen Abstände auf dem Chromosom in Centimorgan (cM).},
    type=fragen
}

\newglossaryentry{wiephyskarte}{
    name={Wie macht man eine physikalische Gen-Karte?},
    description={Man erstellt eine geordnete (BAC-)Klon-Bank, die auf einer Membran repräsentiert ist. Man sucht daraus Klone aus und stattet
    die Enden der Klone mit Sonden aus und markiert sie. Dann erstellt man einen Filter (Membran) mit den Klonen und lässt sie sich hybridisieren und findet so überlappende Klone.
    Basierend auf den Überlappungen erstellt man neue Sonden und wiederholt diesen Prozess mahrmals.
    },
    type=fragen
}


\newglossaryentry{wasistgen}{
    name={Was ist ein Gen},
    description={eine funktionelle DNA-Einheit zur Erzeugung funktioneller RNA-Produkte},
    type=fragen
}
